% SPDX-License-Identifier: CC-BY-4.0
% =====================================================================
% Section 4 of D. Joyner, "A hitch-hiker's guide to the metaplectic trace
% formula for SL(2)", revised for the 2-fold cover (n = 2), November 2026.
%
% Stage 3 of ../revision_log.md. Sources: ../original-1995/cctrace3.tex
% (= ../book-1997/book7c.tex, Sections 4.1-4.2) and the missing cctrac3b.tex
% (= ../book-1997/book7d.tex, Sections 4.3-4.5). Referee items II.5-II.8,
% II.22 (local part), A.1-A.7, the Siegel-domain part of I.10 and the
% Section 4 minor items are fixed here; see the "Stage 3 notes" in
% ../revision_log.md, which also map the 1995 numbers to the numbers here.
%
% This file is \input by the master; it compiles on its own through
% preview.tex. Explicit formulas are checked by checks/sec4_n2_check.py.
% Comments beginning "% [Claude]" are editorial notes by Claude.
% Revised after the second referee report (ChatGPT, November 2026): roadmap
% added, comparisons with the 1995 notes moved to Appendix A.
% =====================================================================

\providecommand{\GA}{\overline{G(\A)}}

\section{The geometric side}\label{sec:geometric}

This section computes the geometric side of the trace formula for
$\GA$, class by class. The general theory is due to
Arthur~\cite{A2,A3,A12} for connected reductive groups and to
Li~\cite{LiI} for covering groups; we specialize it to the double cover of
$G=\Res_{k/\Q}SL(2)$. For this group one inequality between heights
(Lemma~\ref{lem:height-ineq}) makes the reduction theory explicit, with
truncation parameter $T_0=0$. With it, every step of Arthur's coarse
expansion can be checked directly, and each class contributes a
polynomial in $T$ of degree at most one that we write down
(Theorem~\ref{thm:coarse}, Propositions~\ref{prop:hyperbolic},
\ref{prop:elliptic} and~\ref{prop:central}). The new parts are the
corrected weight $2T+v_A$ for hyperbolic classes, the unipotent
contribution of the central classes, given as the finite part of a zeta
integral, and the local distribution $J_A(z,f)$ at a central element
(Lemma~\ref{lem:JA-central}). Section~\ref{sec:fine} states the fine
expansion~\cite{A12,LiI} in the form that Sections~5--7 use.

\paragraph{Roadmap.} The argument has four steps.
\begin{enumerate}[label=\arabic*.]
\item \emph{Classes} (Section~\ref{sec:classes}). $G(\Q)$ is sorted into
  classes $\mathcal O$ by semisimple part: two central classes, the
  hyperbolic classes and the elliptic classes.
\item \emph{Truncation} (Section~\ref{sec:kernels}). The kernel
  $K_G(x,x)=\sum_\gamma f(x^{-1}\gamma x)$ of $R(f)$ is not integrable
  over $G(\Q)\backslash\GA$. The divergence comes from the cusp, the
  region where $H(\delta x)>T$ for some $\delta\in G(\Q)$. There one
  subtracts the kernel $K_B$ of the Borel subgroup; the result is the
  truncated kernel $k^T(x,f)$. The height inequality
  (Lemma~\ref{lem:height-ineq}) says that for $T\ge0$ at most one coset
  $B(\Q)\delta$ is involved at each $x$, so no further reduction theory is
  needed.
\item \emph{Integration, class by class} (Section~\ref{sec:coarse}). For
  an elliptic class nothing is truncated and $J_{\mathcal O}$ is an
  orbital integral. For a hyperbolic class the divergent integral over
  $A(\Q)\backslash A(\A)$ is cut down to an interval of length
  $2T+v_A(y)$, and $J^T_{\mathcal O}$ is a weighted orbital integral. For
  a central class the unipotent elements give a zeta integral, of which
  $J^T_{\mathcal O}$ keeps the finite part at $s=1$ and a term linear in
  $T$.
\item \emph{Localization} (Sections~\ref{sec:local-dist}
  and~\ref{sec:fine}). For $f=f_S\otimes e^S$ the terms are rewritten as
  local distributions on $\overline{G(\Q_S)}$ with global coefficients.
\end{enumerate}
The objects that occur are the following.
\begin{center}
\begin{tabular}{p{0.24\textwidth}p{0.68\textwidth}}
$K_P$, $K_{P,\mathcal O}$ & kernel of $R_P(f)$ and its part on the class
  $\mathcal O$, \eqref{eq:KP}\\
$J_{P,\mathcal O}$ & variant of $K_{P,\mathcal O}$ with the same
  integral, \eqref{eq:JP}\\
$\hat\tau_B(H(\delta x)-T)$ & cutoff: $1$ in the cusp above height $T$,
  $0$ elsewhere\\
$k^T$, $k^T_{\mathcal O}$, $j^T_{\mathcal O}$ & truncated kernels,
  \eqref{eq:kT}\\
$J^T_{\mathcal O}(f)$, $J_{\mathcal O}(f)$ & their integrals, polynomials
  in $T$, and the value at $T=0$\\
$v_A$, $v_{A,v}$ & the weight $-H(x)-H(wx)$ and its local form,
  \eqref{eq:vA}\\
$Z(f,s)$, $\mathrm{FP}_{s=1}$ & zeta integral of the unipotent terms and
  its finite part, \eqref{eq:Zf}\\
$J_G(\gamma,f_S)$, $J_A(\gamma,f_S)$ & local orbital and weighted orbital
  integrals, \eqref{eq:JG}--\eqref{eq:JA}\\
$J_A(z,f_S)$ & the weighted orbital integral at a central element
  (Lemma~\ref{lem:JA-central})\\
$a^M(S,\gamma)$ & global coefficients (Theorem~\ref{thm:fine}); the
  $a^G(S,u_\xi)$ are defined but not evaluated
\end{tabular}
\end{center}

Throughout, $G(\Q)$ is identified with its image in $\GA$ under
Lemma~\ref{lem:global}(a), and $N(\A)$ with $N(\A)\times1$; on $B(\Q)$
the two identifications agree, and $b\mapsto(b,1)$
(Lemma~\ref{lem:global}(b)). Test functions are anti-genuine
(Section~\ref{sec:functions}). We write $w=\begin{pmatrix}0&-1\\1&0
\end{pmatrix}\in G(\Q)$, and $d=[k:\Q]$.

\subsection{Equivalence classes}\label{sec:classes}

Following~\cite[\S2]{A2}, call $\gamma,\gamma'\in G(\Q)$
\emph{equivalent} if their semisimple parts are $G(\Q)$-conjugate, and let
$\mathcal C$ be the set of equivalence classes. A class is
\emph{unramified}~\cite[p.~949]{A2} if it consists of semisimple elements.

\begin{lemma}\label{lem:classes}
The classes $\mathcal O\in\mathcal C$ are of three kinds.
\begin{enumerate}[label=(\alph*)]
\item \emph{Central}: for $z\in\{\pm1\}$, $\mathcal O_z=\{zu\mid
  u\in G(\Q)\ \text{unipotent}\}$. It is not unramified: it contains $z$
  and $z\,n(\nu)$, $\nu\in k^\times$.
\item \emph{Hyperbolic}: the $G(\Q)$-conjugacy class of $\mu=a(t)$,
  $t\in k^\times\smallsetminus\{\pm1\}$. It contains $a(t^{-1})=w\mu w^{-1}$
  and meets $A(\Q)$ in $\{a(t),a(t^{-1})\}$; it is determined by
  $\{t,t^{-1}\}$. The centralizer of $\mu$ is $A$.
\item \emph{Elliptic}: the $G(\Q)$-conjugacy class of a regular
  $\mu\in G(\Q)$ whose characteristic polynomial is irreducible over $k$.
  Its centralizer is a maximal torus $T_\mu=\Res_{k/\Q}T_k$, anisotropic
  over $\Q$, and $\mathcal O\cap A(\Q)=\emptyset$.
\end{enumerate}
Classes of kinds (b) and~(c) are unramified. In all three cases the
centralizer of (the image of) $\mu$ in $\GA$ is the full preimage
$\overline{G_\mu(\A)}$, which is abelian in cases~(b) and~(c).
\end{lemma}

\begin{proof}
An element of $SL(2,k)$ with semisimple part $\pm1$ is $\pm$ a unipotent
element. A regular semisimple element with eigenvalues in $k$ is
$GL(2,k)$-conjugate to $a(t)$, and then $SL(2,k)$-conjugate to it,
because $\diag(y,1)$ centralizes $a(t)$. The centralizers are standard.
For the cover: if $h\in G_\mu(\A)$, lifts of $h$ and $\mu$ commute by
Lemma~\ref{lem:global}(d)--(e) (the argument of
Lemma~\ref{lem:local}(g), applied to the abelian group generated by $h$
and $\mu$); conversely the centralizer in $\GA$ maps into $G_\mu(\A)$.
\end{proof}

For general $n$ the centralizer of a lift of a split $\gamma$ can be
smaller than the preimage of $A$, and then the genuine orbital integral
vanishes. For $n=2$ this does not happen, by the last sentence of the
lemma.

\subsection{Kernels, heights and truncation}\label{sec:kernels}

Let $f\in\mathcal H(\GA)$. We use Arthur's convention: $R$ is the right
regular representation on $L^2(G(\Q)\backslash\GA)_{\mathrm{gen}}$, the
space of $\varphi$ with $\varphi(\gamma x\zeta)=\zeta\varphi(x)$, and
$(R(f)\varphi)(x)=\int_{\GA}f(y)\varphi(xy)\,dy$. The integrand is
invariant under $y\mapsto y\zeta$ because $f$ is anti-genuine. For a
standard parabolic $P=MN_P\in\{B,G\}$ and $\mathcal O\in\mathcal C$ put
\begin{gather}
K_P(x,y)=\sum_{\gamma\in M(\Q)}\int_{N_P(\A)}f(x^{-1}\gamma ny)\,dn,
\qquad
K_{P,\mathcal O}(x,y)=\sum_{\gamma\in M(\Q)\cap\mathcal O}
\int_{N_P(\A)}f(x^{-1}\gamma ny)\,dn,\label{eq:KP}\\
J_{P,\mathcal O}(x,y)=\sum_{\mu\in M(\Q)\cap\mathcal O}\
\sum_{\nu\in N_{P,\mu_s}(\Q)\backslash N_P(\Q)}\ \int_{N_{P,\mu_s}(\A)}
f(x^{-1}\nu^{-1}\mu n\nu y)\,dn,\label{eq:JP}
\end{gather}
where $\mu=\mu_s\mu_u$ is the Jordan decomposition and $N_{P,\mu_s}$ the
centralizer of $\mu_s$ in $N_P$. For $P=G$ these are
$K_G=\sum_{\gamma\in G(\Q)}f(x^{-1}\gamma y)$ and
$K_{G,\mathcal O}=J_{G,\mathcal O}=\sum_{\gamma\in\mathcal O}
f(x^{-1}\gamma y)$.

\begin{lemma}\label{lem:kernels}
The sums~\eqref{eq:KP} and~\eqref{eq:JP} are finite, uniformly for $x,y$
in compact sets. $K_P(x,y)=\sum_{\mathcal O}K_{P,\mathcal O}(x,y)$, and
$K_P$ is the kernel of $R_P(f)$, where $R_P$ is the right regular
representation on $L^2(N_P(\A)M(\Q)\backslash\GA)_{\mathrm{gen}}$.
\end{lemma}

\begin{proof}
$G(\Q)$ is discrete in $G(\A)$ and $f$ has compact support modulo
$\mu_2$, so only finitely many $\gamma$ contribute for $x,y$ in a compact
set; the integrals over $N_P(\A)$ are over compact sets for the same
reason. For upper triangular $b_1,b_2\in B(\A)$ with lower right entries
$d_1,d_2$, \eqref{eq:beta} and $\kappa_v|_B\equiv1$ give
$\sigma_\A(b_1,b_2)=\prod_v(d_2,d_1)_v$, which is $1$ when
$d_1,d_2\in k^\times$. So $b\mapsto(b,1)$ is a homomorphism on
$N(\A)A(\Q)$, and $N_P(\A)M(\Q)$ is a subgroup of $\GA$. For
$\varphi$ in that space,
\[
(R_P(f)\varphi)(x)=\int_{\GA}f(x^{-1}y)\varphi(y)\,dy
=\int_{N_P(\A)M(\Q)\backslash\GA}\ \sum_{\gamma\in M(\Q)}
\int_{N_P(\A)}f(x^{-1}\gamma ny)\,dn\ \varphi(y)\,dy,
\]
using $\varphi(\gamma ny)=\varphi(y)$. The decomposition over
$\mathcal O$ is a regrouping of the finite sum.
\end{proof}

\paragraph{Heights.} For a row vector $u=(c,d)\in k_v^2$ put
$\|u\|_v=\max(|c|_v,|d|_v)$ if $v$ is finite and
$\|u\|_v=(|c|^2+|d|^2)^{[k_v:\R]/2}$ if $v$ is archimedean, and for
$u\in\A_k^2$ put $\|u\|=\prod_v\|u_v\|_v$. Then $\|\lambda u\|_v=
|\lambda|_v\|u\|_v$, and $\|e_2k\|_v=1$ for $k\in K_v$, where
$e_2=(0,1)$. Recall $H=H_B$ from Lemma~\ref{lem:NAK-global}.

\begin{lemma}\label{lem:height-ineq}
\begin{enumerate}[label=(\alph*)]
\item For $x\in\GA$ with image $g\in G(\A)$, $H(x)=-\log\|e_2g\|$.
  The same holds at a single place: $H_v(g)=-\log\|e_2g\|_v$ for
  $g\in SL(2,k_v)$.
\item If $\gamma\in G(\Q)\smallsetminus B(\Q)$, then
  \[
  H(x)+H(\gamma x)\le0\qquad(x\in\GA).
  \]
  Consequently, if $T\ge0$, then for each $x$ there is at most one
  $\delta\in B(\Q)\backslash G(\Q)$ with $H(\delta x)>T$.
\item $H(ax)=H(a)+H(x)$ and $H(wax)=-H(a)+H(wx)$ for
  $a\in\overline{A(\A)}$. Hence
  \begin{equation}\label{eq:vA}
  v_A(x)=-H(x)-H(wx)
  \end{equation}
  is left $\overline{A(\A)}$-invariant, right $\overline{K_0}$-invariant,
  and $v_A(x)\ge0$.
\end{enumerate}
\end{lemma}

\begin{proof}
(a) If $g=n\,a(t)\,k$, then $e_2g=t^{-1}e_2k$, so $\|e_2g\|=|t|^{-1}
=e^{-H(x)}$. The local statement is the same computation.

(b) Let $e_2\gamma=(c,d)$, with $c\ne0$ since $\gamma\notin B(\Q)$. The
matrix with rows $e_2g$ and $e_2\gamma g$ has determinant
$\det\begin{pmatrix}0&1\\c&d\end{pmatrix}\det g=-c$. At each place,
$|\det(u_1;u_2)|_v\le\|u_1\|_v\|u_2\|_v$: this is the ultrametric
inequality at finite places and Hadamard's inequality (squared at complex
places) at archimedean ones. Taking the product over $v$,
$1=|c|_{\A_k}\le\|e_2g\|\,\|e_2\gamma g\|=e^{-H(x)-H(\gamma x)}$ by~(a)
and the product formula. If $\delta_1\ne\delta_2$ in $B(\Q)\backslash
G(\Q)$ both had $H(\delta_ix)>T\ge0$, then (b), applied to $\delta_1x$ and
$\gamma=\delta_2\delta_1^{-1}\notin B(\Q)$, would give $0<2T<
H(\delta_1x)+H(\delta_2x)\le0$.

(c) The first identity is Lemma~\ref{lem:H}(a). For the second,
$wa\,w^{-1}=a^{-1}$, so $H(wax)=H(a^{-1}\cdot wx)=-H(a)+H(wx)$. Right
$K_0$-invariance is Lemma~\ref{lem:H}(b), and $v_A\ge0$ is~(b) with
$\gamma=w$.
\end{proof}

For Lemma~\ref{lem:height-ineq}(b) compare~\cite[Lemma~5.1]{A2}.
Reduction theory~\cite{Go} gives a Siegel set
$\mathfrak S=\{(n,1)\,\tilde a\,\tilde k\mid n\in\omega,\ H(\tilde a)>-c,\
p(\tilde a)\in\Omega\,A_B(\R)^0,\ \tilde k\in\overline{K_0}\}$, with
$\omega\subset N(\A)$ and $\Omega\subset A(\A)^1$ compact and $c\ge0$, such
that $\GA=G(\Q)\mathfrak S$. In the coordinates $x=nak$ the measure
is $\delta_B(a)^{-1}dn\,da\,dk$ by~\eqref{eq:iwasawa}, and $\mathfrak S$ has
finite volume.

\paragraph{Truncation.} Let $\hat\tau_B=\tau_B$ be the characteristic
function of $(0,\infty)$ (Section~\ref{sec:chambers}), $\hat\tau_G=1$, and
$\dim A_G=0$, $\dim A_B=1$. For $T\in\R\cong\mathfrak a$ put
\begin{equation}\label{eq:kT}
k^T(x,f)=K_G(x,x)-\sum_{\delta\in B(\Q)\backslash G(\Q)}
K_B(\delta x,\delta x)\,\hat\tau_B(H(\delta x)-T),
\end{equation}
and define $k^T_{\mathcal O}(x,f)$ and $j^T_{\mathcal O}(x,f)$ in the same
way with $K_{P,\mathcal O}$, respectively $J_{P,\mathcal O}$, in place of
$K_P$. By Lemma~\ref{lem:height-ineq}(b), for $T\ge0$ the sum over
$\delta$ has at most one nonzero term. Let $F^G(x,T)$ be the
characteristic function of $\{x\mid H(\delta x)\le T\ \text{for all}\
\delta\in G(\Q)\}$; its support is compact modulo $G(\Q)$ by reduction
theory~\cite{Go}.

\begin{lemma}\label{lem:partition}
Let $T\ge0$. Then $F^G(x,T)+\sum_{\delta\in B(\Q)\backslash G(\Q)}
\tau_B(H(\delta x)-T)=1$, and
\begin{equation}\label{eq:kT-split}
k^T(x,f)=F^G(x,T)\,K_G(x,x)+\sum_{\delta\in B(\Q)\backslash G(\Q)}
\tau_B(H(\delta x)-T)\bigl(K_G(\delta x,\delta x)-K_B(\delta x,\delta x)\bigr).
\end{equation}
\end{lemma}

\begin{proof}
The first identity says that either $H(\delta x)\le T$ for all $\delta$,
or $H(\delta x)>T$ for exactly one coset, which is
Lemma~\ref{lem:height-ineq}(b). Multiply it by $K_G(x,x)=
K_G(\delta x,\delta x)$ and subtract the $\delta$-sum of~\eqref{eq:kT}.
\end{proof}

In general~\cite[\S6]{A2}, \cite[\S8]{Ar05}, the
partition holds for $T$ sufficiently regular; here $T\ge0$ suffices.

\subsection{The coarse geometric expansion}\label{sec:coarse}

Let $dx$ be a Haar measure on $G(\A)$ with $dx=\delta_B(a)^{-1}dn\,da\,dk$
for $x=nak$ as in~\eqref{eq:iwasawa}, where $dk$ gives $K_0$ volume $1$,
$dn$ gives $N(\Q)\backslash N(\A)$ volume $1$, and $da$ is a Haar measure
on $A(\A)\cong\A_k^\times$, written $d^\times t$ for $a=a(t)$. On covers
$\mu_2$ has volume~$1$. Write $A(\A)=A(\A)^1\times A_B(\R)^0$, where
$A_B(\R)^0=\{a(u)\mid u\in\R_{>0}\}$ ($u$ embedded diagonally in
$k\otimes\R$), let $dX$ be Lebesgue measure on $\R\cong\mathfrak a$
transported to $A_B(\R)^0$ by $H$, and let $da^1$ be the Haar measure on
$A(\A)^1$ with $da=da^1\,dX$. Put
\[
c_k=\mathrm{vol}\bigl(A(\Q)\backslash A(\A)^1\bigr)
=\mathrm{vol}\bigl(k^\times\backslash\A_k^1\bigr).
\]

\begin{theorem}\label{thm:coarse}
Let $f\in\mathcal H(\GA)$ and $T\ge0$.
\begin{enumerate}[label=(\alph*)]
\item Only finitely many $\mathcal O\in\mathcal C$ have
  $k^T_{\mathcal O}(\cdot,f)\not\equiv0$, and $k^T(x,f)=
  \sum_{\mathcal O}k^T_{\mathcal O}(x,f)$.
\item Each $k^T_{\mathcal O}(\cdot,f)$ and each
  $j^T_{\mathcal O}(\cdot,f)$ is absolutely integrable over
  $G(\Q)\backslash\GA$, and
  \[
  J^T_{\mathcal O}(f):=\int_{G(\Q)\backslash\GA}k^T_{\mathcal O}(x,f)\,dx
  =\int_{G(\Q)\backslash\GA}j^T_{\mathcal O}(x,f)\,dx.
  \]
\item Hence $k^T(\cdot,f)$ is absolutely integrable, and
  $J^T(f):=\int k^T(x,f)\,dx=\sum_{\mathcal O}J^T_{\mathcal O}(f)$.
\item Each $J^T_{\mathcal O}(f)$ is a polynomial in $T$ of degree at
  most $1$, given by Propositions~\ref{prop:hyperbolic},
  \ref{prop:elliptic} and~\ref{prop:central}.
\end{enumerate}
\end{theorem}

We put $J_{\mathcal O}(f)=J^0_{\mathcal O}(f)$ and $J(f)=J^0(f)=
\sum_{\mathcal O}J_{\mathcal O}(f)$. In general one defines $J_{\mathcal
O}$ as the value at $T=0$ of the polynomial $J^T_{\mathcal O}$, which is
known only for $T$ sufficiently regular~\cite{A3,Ar05}. Here the integral
itself converges at $T=0$, so the two definitions agree.
Theorem~\ref{thm:coarse} is the case of $\GA$ of Arthur's coarse
expansion~\cite[\S\S7--8]{A2}, \cite{A3}, extended to covers by
Li~\cite{LiI}; for this group we prove it directly.

\begin{proposition}[Hyperbolic classes]\label{prop:hyperbolic}
Let $\mathcal O$ be the class of $\mu=a(t)$, $t\in k^\times
\smallsetminus\{\pm1\}$, and $T\ge0$. Then
\begin{equation}\label{eq:jT-hyp}
j^T_{\mathcal O}(x,f)=\sum_{\eta\in A(\Q)\backslash G(\Q)}
f(x^{-1}\eta^{-1}\mu\eta x)\bigl(1-\hat\tau_B(H(\eta x)-T)
-\hat\tau_B(H(w\eta x)-T)\bigr),
\end{equation}
and
\begin{equation}\label{eq:JT-hyp}
J^T_{\mathcal O}(f)=c_k\int_{\overline{A(\A)}\backslash\GA}
f(y^{-1}\mu y)\,\bigl(2T+v_A(y)\bigr)\,dy,
\end{equation}
with $v_A(y)=-H(y)-H(wy)\ge0$ as in~\eqref{eq:vA}. The integral converges
absolutely.
\end{proposition}

\begin{proof}
Here $\mathcal O\cap A(\Q)=\{\mu,\mu^{-1}\}$ and $N_{\mu_s}=1$, so
\[
J_{B,\mathcal O}(y,y)=\sum_{\mu'\in\{\mu,\mu^{-1}\}}\ \sum_{\nu\in N(\Q)}
f(y^{-1}\nu^{-1}\mu'\nu y).
\] As $\delta$ runs over $B(\Q)\backslash
G(\Q)$ and $\nu$ over $N(\Q)$, $\eta=\nu\delta$ runs over
$A(\Q)\backslash G(\Q)$, and $H(\nu\delta x)=H(\delta x)$. Writing
$\mu^{-1}=w^{-1}\mu w$ and replacing $\eta$ by $w^{-1}\eta$ in the
$\mu^{-1}$-term (note that $H(w^{-1}y)=H(-wy)=H(wy)$)
gives~\eqref{eq:jT-hyp}, since $K_{G,\mathcal O}(x,x)=\sum_{\eta}f(x^{-1}\eta^{-1}\mu\eta x)$.

By Lemma~\ref{lem:height-ineq}(b), for $T\ge0$ the two $\hat\tau_B$ in
\eqref{eq:jT-hyp} are never both $1$, so the bracket is $0$ or $1$.
Unfolding, and writing $x=\tilde a y$ with $\tilde a\in\overline{A(\A)}$,
\[
\int|j^T_{\mathcal O}|\le\int_{\overline{A(\A)}\backslash\GA}
|f(y^{-1}\mu y)|\int_{A(\Q)\backslash\overline{A(\A)}}
\bigl(1-\hat\tau_B(H(\tilde ay)-T)-\hat\tau_B(H(w\tilde ay)-T)\bigr)
\,d\tilde a\,dy,
\]
with equality if $f\ge0$, and we used $f(y^{-1}\tilde a^{-1}\mu\tilde
ay)=f(y^{-1}\mu y)$, because $\overline{A(\A)}$ is abelian
(Lemma~\ref{lem:global}(e)). Put $X=H(\tilde a)$, $h_1=H(y)$, $h_2=
H(wy)$; by Lemma~\ref{lem:height-ineq}(c) the bracket is $1-\mathbf
1[X+h_1>T]-\mathbf 1[-X+h_2>T]$, whose integral over $X\in\R$ is
$2T-h_1-h_2=2T+v_A(y)$. The integral over $A(\Q)\backslash A(\A)^1$
gives $c_k$. This proves~\eqref{eq:JT-hyp} for $\int j^T_{\mathcal O}$
and its absolute convergence: in the coordinates $y=n(z)k$ on
$\overline{A(\A)}\backslash\GA$ (measure $dz\,dk$), $y^{-1}\mu y=
k^{-1}\mu\,n((1-t^{-2})z)k$, so $z$ runs over a compact set and
$v_A(n(z)k)$ is bounded there (Lemma~\ref{lem:weight-local}).

It remains to show $\int k^T_{\mathcal O}=\int j^T_{\mathcal O}$. By
definition
\[
k^T_{\mathcal O}(x)-j^T_{\mathcal O}(x)=\sum_{\delta}
\hat\tau_B(H(\delta x)-T)\,\bigl(J_{B,\mathcal O}-K_{B,\mathcal O}\bigr)
(\delta x,\delta x).
\]
Let $x=(n(u_0),1)\,\tilde a\,\tilde k$ with $p(\tilde a)=a(s)$. For
$\mu'=a(t')$ one has $n(-\nu)\mu'n(\nu)=\mu'n((1-t'^{-2})\nu)$, a
bijection of $N(\Q)$, and $\tilde a^{-1}\mu'(n(z),1)\tilde a=
\mu'(n(s^{-2}z),1)$ (Lemma~\ref{lem:local}(h)). So with
$g_{\mu'}(z)=f(\tilde k^{-1}\mu'(n(z),1)\tilde k)$ and $c=(1-t'^{-2})u_0$,
\[
(J_{B,\mathcal O}-K_{B,\mathcal O})(x,x)=\sum_{\mu'}\Bigl(\sum_{\nu\in k}
g_{\mu'}(s^{-2}(\nu+c))-\int_{\A_k}g_{\mu'}(s^{-2}(z+c))\,dz\Bigr).
\]
By Poisson summation over $k\subset\A_k$~\cite{Ta}, this is
$|s|^2\sum_{\mu'}\sum_{\nu\in k^\times}\hat g_{\mu'}(s^2\nu)\psi(\nu c)$,
which is $O(|s|^{2-M})$ for every $M$, uniformly for $\tilde k\in
\overline{K_0}$ and $s$ in $\Omega\times\R_{\ge1}$, $\Omega\subset
\A_k^1$ compact. Since $\delta_B(a)^{-1}=|s|^{-2}$, the function
$\hat\tau_B(H(x)-T)(J_{B,\mathcal O}-K_{B,\mathcal O})(x,x)$ is
absolutely integrable over $B(\Q)\backslash\GA$, and unfolding is
justified. Integrating first over $u_0\in k\backslash\A_k$ gives $0$,
because $\int_{k\backslash\A_k}\sum_{\nu\in k}g(s^{-2}(\nu+c))\,du_0=
|1-t'^{-2}|_{\A_k}^{-1}\int_{\A_k}g(s^{-2}z)\,dz$ and
$|1-t'^{-2}|_{\A_k}=1$ by the product formula.
\end{proof}

In the form~\eqref{eq:JT-hyp}, $J_{\mathcal O}(f)$ is $c_k$ times Arthur's
weighted orbital integral, $v_A$ being the volume of the convex hull of
$\{-H_B(y),-H_{\bar B}(y)\}$ with $H_{\bar B}(y)=-H(wy)$. For the measure
on $\mathfrak a$ see Remark~\ref{rem:alpha}.

\begin{proposition}[Elliptic classes]\label{prop:elliptic}
Let $\mathcal O$ be elliptic, with $\mu\in\mathcal O$ and $T_\mu=G_\mu$.
Then
$k^T_{\mathcal O}=j^T_{\mathcal O}=K_{G,\mathcal O}(x,x)$ for every $T$,
and
\[
J_{\mathcal O}(f)=J^T_{\mathcal O}(f)=\mathrm{vol}\bigl(T_\mu(\Q)\backslash
T_\mu(\A)\bigr)\int_{\overline{T_\mu(\A)}\backslash\GA}f(y^{-1}\mu y)\,dy,
\]
the integral converging absolutely.
\end{proposition}

\begin{proof}
$\mathcal O\cap A(\Q)=\emptyset$, so the $B$-terms vanish. Unfold
$\sum_{\eta\in T_\mu(\Q)\backslash G(\Q)}$ and use that
$\overline{T_\mu(\A)}$ is abelian, so $f(y^{-1}\tilde t^{-1}\mu\tilde ty)=
f(y^{-1}\mu y)$. The quotient $T_\mu(\Q)\backslash T_\mu(\A)$ is compact
($T_\mu$ is anisotropic, so $T_\mu(\A)^1=T_\mu(\A)$), and the orbit of $\mu$ is closed, so the orbital
integral of the compactly supported $|f|$ converges.
\end{proof}

\paragraph{Central classes.} For $z\in\{\pm1\}$ let $f_z(x)=f(zx)$, with
$z$ identified with its image $(z,1)$ in $\GA$ (Lemma~\ref{lem:global}(a),
$\kappa_\A(\pm1)=1$). It is central, so $f_z\in\mathcal H(\GA)$, and at
almost all $v$ its local component is again $e_v$, because $(-1,1)\in
K_v^*$. For $f\in\mathcal H(\GA)$ define a Schwartz--Bruhat function on
$\A_k$ by
\begin{equation}\label{eq:Phi}
\Phi_f(y)=\int_{\overline{K_0}}f\bigl(\tilde k^{-1}(n(y),1)\tilde k\bigr)
\,d\tilde k,
\end{equation}
let $\hat\Phi_f$ be its Fourier transform with respect to $\psi$, and for
$\mathrm{Re}\,s>1$ put
\begin{equation}\label{eq:Zf}
Z(f,s)=\int_{k^\times\backslash\A_k^\times}|t|^{-2s}\sum_{\nu\in k^\times}
\Phi_f(t^{-2}\nu)\,d^\times t
=\frac12\sum_{\xi\in k^\times/k^{\times2}}\ \int_{\A_k^\times}
\Phi_f(\xi t^2)\,|t|^{2s}\,d^\times t .
\end{equation}
For $f=\bigotimes_vf_v$ the last integral is the product of the local
integrals $Z_v(\Phi_{f,v},\xi,s)=\int_{k_v^\times}\Phi_{f,v}(\xi t^2)
|t|_v^{2s}d^\times t$, which are, at $s=1$, twice the local unipotent orbital
integrals of $f_v$ at $n(\xi)$ (Lemma~\ref{lem:JA-central}(c)).

\begin{proposition}[Central classes]\label{prop:central}
Let $z\in\{\pm1\}$. Then $k^T_{\mathcal O_z}=j^T_{\mathcal O_z}$, and
$J^T_{\mathcal O_z}(f)=J^T_{\mathrm{unip}}(f_z)$, where for every $T\in\R$
\begin{equation}\label{eq:Junip}
J^T_{\mathrm{unip}}(f)=\mathrm{vol}\bigl(G(\Q)\backslash\GA\bigr)f(1)
+c_k\,\hat\Phi_f(0)\,T+\mathop{\mathrm{FP}}_{s=1}Z(f,s).
\end{equation}
Here $Z(f,s)$, absolutely convergent for $\mathrm{Re}\,s>1$, extends to a
meromorphic function on $\C$ whose only possible poles are simple ones at
$s=1$ and $s=0$; the residue at $s=1$ is $c_k\hat\Phi_f(0)/2$, and
$\mathrm{FP}_{s=1}$ denotes the constant term of the Laurent expansion.
In particular
\[
J_{\mathcal O_z}(f)=J_{\mathrm{unip}}(f_z),\qquad
J_{\mathrm{unip}}(f)=\mathrm{vol}\bigl(G(\Q)\backslash\GA\bigr)f(1)
+\mathop{\mathrm{FP}}_{s=1}Z(f,s).
\]
\end{proposition}

\begin{proof}
Since $z$ is central, $K_{P,\mathcal O_z}(x,y)$ for $f$ equals
$K_{P,\mathcal O_1}(x,y)$ for $f_z$, and likewise for $J_P$; and
$J_{P,\mathcal O_1}=K_{P,\mathcal O_1}$ because $\mu_s=1$ and
$N_{P,1}=N_P$. So it suffices to treat $z=1$. The nontrivial unipotent
elements of $G(\Q)$ are the $\delta^{-1}n(\nu)\delta$,
$\delta\in B(\Q)\backslash G(\Q)$, $\nu\in k^\times$, each written once,
so
\[
k^T_{\mathcal O_1}(x)=f(1)+\sum_{\delta}\Bigl(F(\delta x)
-\hat\tau_B(H(\delta x)-T)\,C(\delta x)\Bigr),
\]
where $F(x)=\sum_{\nu\in k^\times}f(x^{-1}n(\nu)x)$ and
$C(x)=\int_{N(\A)}f(x^{-1}nx)\,dn$.
For $x=(n_0,1)\tilde a\tilde k$ with $p(\tilde a)=a(t)$ we have
$x^{-1}(n(\nu),1)x=\tilde k^{-1}(n(t^{-2}\nu),1)\tilde k$, so $F$ and $C$
depend only on $t$ and $\tilde k$. With $g_{\tilde k}(y)=
f(\tilde k^{-1}(n(y),1)\tilde k)$, Poisson summation gives
\[
\delta_B(a)^{-1}\bigl(F-C\bigr)(x)=\sum_{\nu\in k^\times}
\hat g_{\tilde k}(t^2\nu)-|t|^{-2}g_{\tilde k}(0),
\]
which is integrable over $|t|>e^T$, uniformly in $\tilde k$; and
$F(x)=0$ once $|t|$ is small, because $g_{\tilde k}$ has compact support
uniformly in $\tilde k$ and $k$ is discrete in $\A_k$. Hence the
three pieces $f(1)$, $\sum_\delta(1-\hat\tau_B)F$ and
$\sum_\delta\hat\tau_B(F-C)$ of $k^T_{\mathcal O_1}$ are absolutely
integrable, and unfolding each over $B(\Q)\backslash\GA=N(\Q)\backslash
N(\A)\times A(\Q)\backslash A(\A)\times\overline{K_0}$ gives, with
$G(t)=|t|^{-2}\sum_{\nu\in k^\times}\Phi_f(t^{-2}\nu)$ and
$\delta_B(a)^{-1}C=\hat\Phi_f(0)$ after the $\tilde k$-integration,
\[
J^T_{\mathrm{unip}}(f)=\mathrm{vol}\bigl(G(\Q)\backslash\GA\bigr)f(1)
+\int_{|t|\le e^T}G(t)\,d^\times t+\int_{|t|>e^T}\bigl(G(t)-\hat\Phi_f(0)
\bigr)\,d^\times t.
\]
Now $Z(f,s)=\int G(t)|t|^{2-2s}d^\times t$. Split it at $|t|=e^T$ and
write $G=(G-\hat\Phi_f(0))+\hat\Phi_f(0)$ on $|t|>e^T$. Since
$\int_{|t|>e^T}|t|^{2-2s}d^\times t=c_k\,e^{(2-2s)T}/(2s-2)$,
\[
Z(f,s)=\int_{|t|\le e^T}G|t|^{2-2s}+\int_{|t|>e^T}\bigl(G-\hat\Phi_f(0)
\bigr)|t|^{2-2s}+c_k\hat\Phi_f(0)\frac{e^{(2-2s)T}}{2s-2}.
\]
The first integral is entire in $s$ (compact range of $|t|$). By the
Poisson formula above, $G-\hat\Phi_f(0)=\sum_{\nu\ne0}\hat\Phi_f(t^2\nu)
-|t|^{-2}\Phi_f(0)$, so the second integral is an entire function minus
$c_k\Phi_f(0)e^{-2sT}/(2s)$. This gives the continuation and the poles.
Expanding $e^{(2-2s)T}/(2s-2)=1/(2(s-1))-T+O(s-1)$ gives~\eqref{eq:Junip}.
The second form of~\eqref{eq:Zf} follows by writing $k^\times$ as the
union of the cosets $\xi k^{\times2}$ and unfolding
$\sum_{c\in k^\times/\{\pm1\}}$ against $k^\times\backslash\A_k^\times$,
using $|ct|=|t|$.
\end{proof}

The linear term of~\eqref{eq:Junip} agrees with the fine expansion: the
term $\tfrac12a^A(S,1)J^T_A(1,f)$ there has weight $v_A+2T$, and
$a^A(S,1)=c_k$, $J_G$ of the class induced from $1$ tends to
$\hat\Phi_f(0)$ (Lemma~\ref{lem:JA-central}).

\begin{example}\label{ex:unip-Q}
Let $k=\Q$ and $\Phi=\Phi_\infty\otimes\mathbf 1_{\hat\Z}$ with
$\Phi_\infty(y)=e^{-\pi y^2}$, so $\hat\Phi(0)=1$, and take $d^\times t$
with $\mathrm{vol}(\hat\Z^\times)=1$, so $c_\Q=1$. Representing
$\Q^\times\backslash\A^\times$ by $\R_{>0}\times\hat\Z^\times$,
$Z(s)=\int_0^\infty u^{-2s}\sum_{\nu\ne0}e^{-\pi\nu^2/u^4}\,du/u=
\zeta(s)M(s)$ with $M(s)=\int_0^\infty e^{-\pi w^2}w^{s-1}dw=
\tfrac12\pi^{-s/2}\Gamma(s/2)$. The residue at $s=1$ is $M(1)=\tfrac12=
c_\Q\hat\Phi(0)/2$, and
\[
\mathop{\mathrm{FP}}_{s=1}Z(s)=\gamma_EM(1)+M'(1)
=\frac{\gamma_E}4-\frac{\log4\pi}4\approx-0.48845 .
\]
In the second form of~\eqref{eq:Zf}, $\xi$ runs over $\pm m$ with $m$
squarefree, each contributing $m^{-s}M(s)\zeta(2s)$, and
$\sum_m m^{-s}=\zeta(s)/\zeta(2s)$. The check script computes the
truncated integrals of the proof of Proposition~\ref{prop:central} at
$T=-0.3,0,0.5,1.2$ and finds that
$J^T_{\mathrm{unip}}-\mathrm{vol}\cdot f(1)-T$ equals this constant to
$14$ digits.
\end{example}

\begin{proof}[Proof of Theorem~\ref{thm:coarse}]
(a) $k^T(x)=\sum_{\mathcal O}k^T_{\mathcal O}(x)$ by
Lemma~\ref{lem:kernels}. If $k^T_{\mathcal O}\not\equiv0$, then some
element of $\mathcal O$, or some $\mu n$ with $\mu\in\mathcal O\cap A(\Q)$,
is $G(\A)$-conjugate into the support of $f$; its trace lies in $k$ and
in a compact subset of $\A_k$, so in a finite set. This leaves the two
central classes, finitely many hyperbolic classes, and the elliptic
classes with finitely many characteristic polynomials. The
$SL(2,k)$-classes with a given regular elliptic characteristic polynomial,
splitting field $E$, form a torsor under $k^\times/N_{E/k}E^\times$,
which injects into $\bigoplus_vk_v^\times/N E_v^\times$ by the Hasse
norm theorem~\cite{CF}. Let $v$ be a finite place outside a finite set,
where $f_v=e_v$, $E_v/k_v$ is unramified and $\mu$ is a unit with
$|D(\mu)|_v=1$. If $v$ splits in $E$, the local class is unique. If $v$
is inert, an element $\gamma\in K_v$ with the characteristic polynomial
of $\mu$ has irreducible reduction, so it fixes no edge of the
Bruhat--Tits tree and exactly one vertex, that of $K_v$. If $\gamma'=
m\gamma m^{-1}\in K_v$ with $m\in GL(2,k_v)$, then $m$ fixes that vertex,
so $m\in k_v^\times GL(2,\OO_v)$ and $\det m\in k_v^{\times2}\OO_v^\times
=N E_v^\times$; correcting $m$ by an element of the centralizer $E_v^\times$
makes it lie in $SL(2,k_v)$. So one local class meets $K_v$. Hence only
finitely many local classes, and so finitely many global classes, meet the
support of $f$.

(b) and (d) are Propositions~\ref{prop:hyperbolic}, \ref{prop:elliptic}
and~\ref{prop:central}. (c) follows from (a) and (b).
\end{proof}

\subsection{Local distributions}\label{sec:local-dist}

\paragraph{Test functions.} Let $S$ be a finite set of places of $\Q$
containing $\infty$ and $2$, and let $S_k$ be the set of places of $k$
above $S$. Let $\overline{G(\Q_S)}$ be the cover of $G(\Q_S)=
\prod_{v\in S_k}SL(2,k_v)$ defined by $\prod_{v\in S_k}\sigma_v$, the
quotient of $\prod_{v\in S_k}\overline{SL(2,k_v)}$ as in
Section~\ref{sec:local}. Let $\mathcal H(\overline{G(\Q_S)})$ be the span
of the products $\bigotimes_{v\in S_k}f_v$, $f_v\in
C_c^\infty(\overline{SL(2,k_v)})_-$, with $f_v$ bi-$\overline{K_v}$-finite
at archimedean $v$; here $C_c^\infty(\cdot)_-$ is the space of
Section~\ref{sec:functions}, defined on the topological cover. Let
$\mathcal H_N(\overline{G(\Q_S)})_\Gamma$, for $N>0$ and a finite set $\Gamma$ of irreducible
representations of $\overline{K_S}$, be the subspace of functions
supported on $\{\|g\|\le e^N\}$ whose left and right
$\overline{K_S}$-translates span representations with constituents in
$\Gamma$. With the seminorms $\sup|Df|$, $D$ a differential operator on
$\overline{G(\R)}$, it is a Fr\'echet space, and
$\mathcal H(\overline{G(\Q_S)})$ is the union over $N$ and $\Gamma$. For
$f_S\in\mathcal H(\overline{G(\Q_S)})$ put $f=f_S\otimes e^S$, $e^S=
\bigotimes_{v\notin S_k}e_v$. For $\gamma\in G(\Q)$ let
$\gamma_S=(\gamma,\kappa_\A(\gamma))\in\overline{G(\Q_S)}$; since
$e_v((k,\zeta)_\sigma)=\zeta$ for $k\in K_v$, $f(\gamma)=f_S(\gamma_S)$
when $\gamma\in K_v$ for all $v\notin S_k$.

Arthur's spaces of almost compactly supported and moderate
functions~\cite{A9} coincide with $\mathcal H$ for $G$, because $A_G=1$.
For $A$ they do not: $\mathfrak a_A=\mathfrak a\ne0$, and a function that
is compactly supported on each fiber of $H$ need not be compactly
supported. This matters in Section~\ref{sec:IAgamma}.

\paragraph{Orbital integrals.} At each $v\in S_k$ we use the local
measures of Section~\ref{sec:measures-local}, and on $G(\Q_S)$ their
product. For $\gamma=a(t)\in A(\Q_S)$ with $t_v^2\ne1$ for all $v$, and
$\gamma$ also denoting a lift in $\overline{A(\Q_S)}$, let
$D(\gamma)=\prod_v(t_v-t_v^{-1})^2$, and
\begin{align}
J_G(\gamma,f_S)&=|D(\gamma)|^{1/2}\int_{A(\Q_S)\backslash G(\Q_S)}
f_S(x^{-1}\gamma x)\,dx,\label{eq:JG}\\
J_A(\gamma,f_S)&=|D(\gamma)|^{1/2}\int_{A(\Q_S)\backslash G(\Q_S)}
f_S(x^{-1}\gamma x)\,v_A(x)\,dx,\qquad
v_A(x)=\sum_{v\in S_k}v_{A,v}(x_v),\label{eq:JA}
\end{align}
where $v_{A,v}(x)=-H_v(x)-H_v(wx)$ is the local weight. For
$\gamma=\mu_S$ regular elliptic, $J_G(\gamma,f_S)$ is defined
by~\eqref{eq:JG} with the centralizer $T_\mu(\Q_S)$ in place of
$A(\Q_S)$.
For a unipotent $u\in G(\Q_S)$ and $z\in\{\pm1\}$, $J_G(zu,f_S)=
\int_{G_u(\Q_S)\backslash G(\Q_S)}f_S(x^{-1}zux)\,dx$. These integrals
converge absolutely, and the quotient measures are those of~\eqref{eq:iwasawa}:
on $A(\Q_S)\backslash G(\Q_S)$, $dx=dn\,dk$ for $x=nk$.

\begin{lemma}\label{lem:weight-local}
Let $v$ be a place of $k$. Then $v_{A,v}\ge0$, it is left
$A(k_v)$-invariant and right $K_v$-invariant, and
\[
v_{A,v}(n(y)k)=\log\|(1,y)\|_v=\begin{cases}
\log\max(1,|y|_v),&v\ \text{finite},\\
\tfrac12\log(1+y^2),&v\ \text{real},\\
\log(1+|y|^2),&v\ \text{complex}.\end{cases}
\]
For $x\in G(\A)$, $v_A(x)=\sum_vv_{A,v}(x_v)$, a finite sum.
\end{lemma}

\begin{proof}
As for Lemma~\ref{lem:height-ineq}, at one place: $|\det(e_2x;e_2wx)|_v
=|-1|_v=1$ gives $v_{A,v}\ge0$, and $H_v(n(y))=0$, $e_2wn(y)=(1,y)$. The
global identity is $H=\sum_vH_v$.
\end{proof}

\begin{lemma}\label{lem:JA-central}
Let $f_S=\bigotimes_{v\in S_k}f_v$, $z\in\{\pm1\}$, and let
$\Phi_v(y)=\int_{\overline{K_v}}f_v(\tilde k^{-1}z(n(y),1)\tilde k)\,
d\tilde k$, a Schwartz--Bruhat function on $k_v$, with $\hat\Phi_v(0)=
\int_{k_v}\Phi_v$. For $t\in k_S^\times$ near $1$ let $a_t=(a(t),1)z$
and
\[
r(t)=\sum_{v\in S_k}\log\bigl|t_v^2-t_v^{-2}\bigr|_v .
\]
\begin{enumerate}[label=(\alph*)]
\item $J_G(a_t,f_S)\to\prod_v\hat\Phi_v(0)$ as $t\to1$.
\item The limit
  \[
  J_A(z,f_S):=\lim_{t\to1}\bigl(J_A(a_t,f_S)+r(t)\,J_G(a_t,f_S)\bigr)
  =\sum_{v\in S_k}\Bigl(\int_{k_v}\Phi_v(y)\log|2y|_v\,dy\Bigr)
  \prod_{w\in S_k\smallsetminus\{v\}}\hat\Phi_w(0)
  \]
  exists. Without the term $r(t)J_G$, $J_A(a_t,f_S)=-\log|1-t^{-2}|_S
  \prod_v\hat\Phi_v(0)+O(1)$, which diverges unless
  $\prod_v\hat\Phi_v(0)=0$.
\item For $\xi\in k_S^\times$, $J_G(z\,n(\xi),f_S)=\prod_{v\in S_k}
  \tfrac12\int_{k_v^\times}\Phi_v(\xi t^2)|t|_v^2\,d^\times t$.
\end{enumerate}
\end{lemma}

\begin{proof}
Everything factors over $v\in S_k$, $J_A$ by the Leibniz rule
$J_A=\sum_vJ_{A,v}\prod_{w\ne v}J_{G,w}$, and $r(t)=\sum_vr_v(t_v)$; so
it suffices to treat one place. For $x=n(y)k$,
$x^{-1}a_tx=\tilde k^{-1}a_t(n(\varepsilon y),1)\tilde k$ with
$\varepsilon=1-t^{-2}$, and $|D(a(t))|^{1/2}=|t|\,|\varepsilon|$. Put
$y'=\varepsilon y$ and $h_t(y')=\int f_v(\tilde k^{-1}a_t(n(y'),1)
\tilde k)\,d\tilde k$. By Lemma~\ref{lem:weight-local} and homogeneity of
$\|\cdot\|_v$,
\[
J_{G,v}=|t|\int h_t(y')\,dy',\qquad
J_{A,v}=|t|\int h_t(y')\bigl(\log\|(\varepsilon,y')\|_v-\log|\varepsilon|_v
\bigr)\,dy'.
\]
This gives~(a), since $h_t\to\Phi_v$ with uniformly compact support, and
the asymptotics in~(b). Since $t^2-t^{-2}=t^2\varepsilon(1+t^{-2})$,
\[
J_{A,v}+r_vJ_{G,v}=|t|\int h_t(y')\bigl(\log\|(\varepsilon,y')\|_v+
\log|t^2(1+t^{-2})|_v\bigr)\,dy'\to\int\Phi_v(y)\bigl(\log|y|_v+
\log|2|_v\bigr)\,dy
\]
by dominated convergence ($\log|y|_v$ is locally integrable). For~(c),
$G_{n(\xi)}=\{\pm1\}N$, so the integral is over $\{\pm1\}\backslash
A(k_v)\times K_v$ with measure $\tfrac12\delta_B(a)^{-1}da\,dk$, and
$a(t)^{-1}n(\xi)a(t)=n(t^{-2}\xi)$; substitute $t\mapsto t^{-1}$.
\end{proof}

Lemma~\ref{lem:JA-central}(b) is Arthur's definition~\cite[(8.1)]{A7} of
$J_M(\gamma,f)$ at a non-regular point, as the limit of
$\sum_{L\in\mathcal L(M)}r_M^L(\gamma,a)J_L(a\gamma,f)$, specialized to
$M=A$, $\gamma=z$: $r^A_A=1$, and $r^G_A(1,a)$ is the function $r(t)$
above (Remark~\ref{rem:alpha}). The term $L=G$ is present because the
class induced from $1\in A$ to $G$ is the regular unipotent class; without
it the limit diverges, as~(b) shows.

\subsection{The fine expansion}\label{sec:fine}

For $S$ as above let $\mathcal U_S$ be the set of $(G,S)$-equivalence
classes of unipotent elements of $G(\Q)$, that is, the classes
$\{1\}$ and $u_\xi=n(\xi)$, $\xi\in k_S^\times/k_S^{\times2}$ (each
meets $G(\Q)$, by weak approximation). For $v\notin S_k$ we use Haar
measures giving $K_v$, $A(\OO_v)$ and $T_\mu(\OO_v)$ volume $1$, and global
measures that are products of these with the measures at $S_k$; the
volumes below are taken with respect to them.

\begin{theorem}[Arthur~{\cite[Thm.~8.1]{A12}}; Li~\cite{LiI} for
covers]\label{thm:fine}
Let $f_S\in\mathcal H(\overline{G(\Q_S)})$ and $f=f_S\otimes e^S$. If
$S$ is large enough (depending on the support of $f_S$), then:
\begin{enumerate}[label=(\alph*)]
\item for $z\in\{\pm1\}$,
  \[
  J_{\mathcal O_z}(f)=a^G(S,1)f_S(z)+\sum_{\xi\in k_S^\times/k_S^{\times2}}
  a^G(S,u_\xi)\,J_G(z\,u_\xi,f_S)+\tfrac12\,a^A(S,1)\,J_A(z,f_S),
  \]
  with $a^G(S,1)=\mathrm{vol}(G(\Q)\backslash\GA)$, $a^A(S,1)=c_k$, and
  $J_A(z,f_S)$ as in Lemma~\ref{lem:JA-central}(b); the coefficients
  $a^G(S,u_\xi)$ are those of~\cite{A12,LiI} and are not evaluated in this
  paper;
\item for $\mathcal O$ elliptic with $\mu\in\mathcal O$,
  $J_{\mathcal O}(f)=a^G(S,\mu)J_G(\mu_S,f_S)$ with
  $a^G(S,\mu)=\mathrm{vol}(T_\mu(\Q)\backslash T_\mu(\A))$;
\item for $\mathcal O$ hyperbolic with $\mu=a(t)\in\mathcal O$,
  $J_{\mathcal O}(f)=a^A(S,\mu)J_A(\mu_S,f_S)$ with $a^A(S,\mu)=c_k$.
\end{enumerate}
\end{theorem}

In Arthur's notation the right-hand sides are
$\sum_{M\in\{A,G\}}|W^M_0||W^G_0|^{-1}\sum_{\gamma}a^M(S,\gamma)
J_M(\gamma,f_S)$, with $\gamma$ running over the $(M,S)$-classes in
$M(\Q)\cap\mathcal O$; in~(c) the factor $\tfrac12$ cancels against the
two elements $\mu^{\pm1}$, since $J_A(\mu^{-1},f_S)=J_A(\mu,f_S)$. (Substitute
$x=wy$, and use $v_A(wy)=-H(wy)-H(-y)=v_A(y)$.) Parts~(b) and~(c) follow from
Propositions~\ref{prop:elliptic} and~\ref{prop:hyperbolic}: for $S$
large, $|D(\mu)|_S=|D(\mu)|_{\A_k}=1$, and at $v\notin S_k$ the orbital
integrals of $e_v$ equal $1$ while $v_{A,v}=0$ on their support. For the
same reason $f(\gamma)=f_S(\gamma_S)$, which is why the lift $\mu_S$
carries the global factor $\kappa_\A(\mu)$. In~(a) the coefficient of
$f_S(z)$ and the linear term agree with Proposition~\ref{prop:central}.
The coefficients $a^G(S,u_\xi)$ exist by the theorem and are determined
by~(a), but we do not compute them (Remark~\ref{rem:sec4-status}).
Wherever they occur below, in Theorems~\ref{thm:tf} and~\ref{thm:itf}, they
are defined, not evaluated.

The coefficients are defined with centralizers in the cover, which for
$n=2$ are the full preimages (Lemma~\ref{lem:classes}), and the sums run
over $(G,S)$-classes, not over the rational class $\mathcal O$.

\begin{remarkrev}\label{rem:alpha}
\emph{Normalizations on $\mathfrak a$.} Identify $\mathfrak a$ with $\R$ by
$X\mapsto X(\chi_1)$, $\chi_1=N_{k/\Q}\circ t$, as in
Section~\ref{sec:chambers}. The coroot is $\alpha^\vee\colon s\mapsto
a(s)$, $s\in\Q^\times$, so $\alpha^\vee=\langle\chi_1,\alpha^\vee\rangle=
d$, and the simple root $\alpha$ of $(B,A_B)$, $a(s)\mapsto s^2$, is
$\alpha=(2/d)\chi_1$, i.e.\ $\alpha(X)=2X/d$. The modulus is $\delta_B=e^{2H}$ for every $k$. We use
Lebesgue measure $dX$ on $\mathfrak a$, for which $\mathfrak
a/\Z\alpha^\vee$ has volume $d$; Arthur normalizes it to volume $1$. With
Arthur's normalization, $\theta_B(\lambda)=\lambda(\alpha^\vee)$, the
weight $v_A$ and the function $r_A^G$ are divided by $d$, and
$\mathrm{vol}(A(\Q)\backslash A(\A)^1)$ is multiplied by $d$; the products
in~\eqref{eq:JT-hyp} and Theorem~\ref{thm:fine} do not change. With
$\theta_P(\lambda)$ in our normalization, Arthur's
$r^G_A(1,a)=\lim_{\lambda\to0}\sum_{P\in\mathcal P(A)}r_P(\lambda,1,a)
\theta_P(\lambda)^{-1}$~\cite[\S3]{A7} with $r_P(\lambda,1,a)=
|a^\beta-a^{-\beta}|^{\lambda(\beta^\vee)/2}$ evaluates to $r(t)$ for
$a=a(t)$, $t\in\Q_S^\times$. We have checked this only for the formula
of~\cite{A7} as quoted in \texttt{book7d.tex}, not against the general
definition of the exponents $\rho(\beta,u)$ there. The independent check is
that the limit in Lemma~\ref{lem:JA-central}(b) exists.
\end{remarkrev}

\subsection{Status of the results}\label{sec:status4}

\begin{remarkrev}\label{rem:sec4-status}
\emph{What rests on citation.} Theorem~\ref{thm:coarse} and
Propositions~\ref{prop:hyperbolic}--\ref{prop:central} are proved here.
Theorem~\ref{thm:fine}(a) is taken from~\cite{A12,LiI}: the coefficients
$a^G(S,u_\xi)$ are defined by it, and we have not derived~(a) from
Proposition~\ref{prop:central}. Doing so amounts to computing
$a^G(S,u_\xi)$ from the Laurent expansion at $s=1$ of the second form
of~\eqref{eq:Zf}, grouped by the image of $\xi$ in
$k_S^\times/k_S^{\times2}$. Each fiber is infinite, and its contribution
is a Dirichlet series in $s$, so a regularization is needed. This is
left to the revision of Section~7, where the same divergence occurs. The unramified computations used for (b)--(c) of
Theorem~\ref{thm:fine} are standard and not written out. For an
elliptic $\mu$ they use the uniqueness of the local class meeting $K_v$
from the proof of Theorem~\ref{thm:coarse}(a).
\end{remarkrev}

\begin{remarkrev}\label{rem:verified-sec4}
The script \texttt{checks/sec4\_n2\_check.py} verifies, for $k=\Q$: that
the Iwasawa height computed from an explicit decomposition equals
$-\log\|e_2g\|_v$ at $v=\infty,2,3,5,7$; the inequality
$H(x)+H(\gamma x)\le0$ on random adelic $x$ and $\gamma\in
SL(2,\Q)\smallsetminus B(\Q)$, and $H(\gamma x)=H(x)$ for $\gamma\in
B(\Q)$; Lemma~\ref{lem:height-ineq}(c) and Lemma~\ref{lem:weight-local};
the weight integral in the proof of Proposition~\ref{prop:hyperbolic};
the conjugation identities used there (symbolically);
Lemma~\ref{lem:JA-central}(b) at $v=\R$ (quadrature) and $v=\Q_3$
(exactly), including the divergence without $r(t)$; the symmetry
$J_A(\mu^{-1})=J_A(\mu)$ at $\R$; and Example~\ref{ex:unip-Q}.
\end{remarkrev}
